% Supplement: proofs of statements in the main text.
% Sources: draft/bie-enclosure/s-modelproblem.tex, a-cellrep.tex, a-trace.tex,
% a-reliable.tex (clean version).
\section{Proofs}\label{sm:proofs}

\subsection{Proof of \texorpdfstring{\cref{prop:ess}}{Proposition 2.1}}\label{sm:proof-ess}
The second identity is the Floquet--Bloch decomposition in $x$: $A_s(\beta)$ is
unitarily equivalent to the direct integral of the fibers $A_s(\beta,\theta)$,
each band is continuous (\cref{lem:lipschitz}), and
$\inf_\theta\lambda_{s,n}(\theta)\to\infty$, so the union of bands is locally
finite and closed.

($\subset$) Let $\lambda\notin\sigma(A_-)\cup\sigma(A_+)$ be real and set
$R_\pm=(A_\pm-\lambda)^{-1}$, $R_0=(A+1)^{-1}$. Take real smooth functions
$\chi_-,\chi_0,\chi_+$ of $x$ alone with $\chi_-^2+\chi_0^2+\chi_+^2=1$,
$\chi_0$ compactly supported, and $\chi_\pm$ together with their derivatives
supported in the regions where $\rho=\rho_\pm$. The weighted norms are
equivalent, so $Q=\chi_-R_-\chi_-+\chi_0R_0\chi_0+\chi_+R_+\chi_+$ maps
$\mathcal H$ boundedly into $D(A)$. On the support of $\chi_\pm$ we have
$(A-\lambda)\chi_\pm R_\pm\chi_\pm=\chi_\pm^2+[A,\chi_\pm]R_\pm\chi_\pm$, and
similarly for $\chi_0$; hence $(A-\lambda)Q=I+K$, where $K$ is a finite sum of
terms $[A,\chi]R\chi$, $[A,\chi]v=-\rho^{-1}(2\chi'\partial_xv+\chi''v)$, and of
$-(\lambda+1)\chi_0R_0\chi_0$. Each term factors through a compact embedding
$H^2\hookrightarrow H^1$ or $H^2\hookrightarrow L^2$ on a bounded $x$-window, so
$K$ is compact. Thus $A-\lambda$ has a bounded right parametrix, closed range of
finite codimension, and by self-adjointness a finite-dimensional kernel: it is
Fredholm, and $\lambda\notin\sigma_{\mathrm{ess}}(A)$.

($\supset$) Let $\lambda=\lambda_{s,n}(\theta)$ with a nonzero Bloch function
$w$. For a nonzero $\eta\in C_c^\infty(\RR)$ set
$u_N(x,y)=c_N\eta((x-X_N)/(NL_s))w(x,y)$, where $X_N$ moves the support into
lead $s$ and to infinity, and $c_N=O(N^{-1/2})$ normalizes $u_N$. On the
support $\rho=\rho_s$, so the residual comes only from $\eta'$ and $\eta''$:
$\|(A-\lambda)u_N\|_{\mathcal H}=O(N^{-1})$ and $u_N\rightharpoonup0$. By Weyl's
criterion $\lambda\in\sigma_{\mathrm{ess}}(A)$.

\subsection{Proof of \texorpdfstring{\cref{prop:discrete}}{Proposition 2.3}}\label{sm:proof-discrete}
By \cref{prop:ess} there is no essential spectrum in $G_\beta$, and a
nonessential spectral point of a self-adjoint operator is an isolated
eigenvalue of finite multiplicity. For identical leads the decay is
\cite[Thm.~3.5]{Fliss2013}; we treat the right lead. Let $Au=\lambda u$ with
$\lambda\in G_\beta$ and $g=\operatorname{dist}(\lambda,\sigma(A_+))>0$. With a
cutoff $\eta(x)$ that rises from $0$ to $1$ inside the right lead, $w=\eta u$
solves $(A_+-\lambda)w=f$ with $f$ supported in a bounded $x$-window. Take a
bounded weight $h(x)$ that vanishes on that window, has $|h'|\le\alpha$ and
$|h''|\le\alpha/\ell$, grows with slope $\alpha$ to the right and is flattened
far away (with $\ell$ independent of the flattening point); then $\ee^hf=f$ and
$\ee^h(A_+-\lambda)\ee^{-h}=A_+-\lambda+\rho_+^{-1}(2h'\partial_x+h''-h'^2)$.
By the spectral theorem $\|(A_+-\lambda)^{-1}\|\le g^{-1}$ and
$\|A_+^{1/2}(A_+-\lambda)^{-1}\|^2\le g^{-1}+\lambda g^{-2}$, so the
perturbation has relative norm at most
\[
  q(\alpha)=\frac{2\alpha(\lambda+g)^{1/2}}{\rin^{1/2}g}
  +\frac{\alpha^2+\alpha/\ell}{\rin g}.
\]
For $q(\alpha)<1$ a Neumann series gives
$\|\ee^hw\|_{\mathcal H_+}\le\|\ee^hf\|_{\mathcal H_+}/(g(1-q))
=\|f\|_{\mathcal H_+}/(g(1-q))$, uniformly in the flattening point. The same
resolvent bound applied to
$(A_+-\lambda)(\ee^hw)=f-\rho_+^{-1}(2h'\partial_x+h''-h'^2)\ee^hw$ bounds the
graph norm of $\ee^hw$ uniformly; letting the flattening point go to infinity
gives $\ee^{\alpha x}u\in H^2$ far to the right. The left lead is the same.

\subsection{Proof of \texorpdfstring{\cref{thm:complement}}{Theorem 2.6}}\label{sm:proof-complement}
Let $b[u]=\int\rho_s|u|^2$ on the cell, so that the eigenvalues
$\lambda_{s,j}(\theta)$ are those of the pencil $(a_\theta,b)$ on the gauged form
domain. Let $P$ be the $L^2$-orthogonal projection onto
$\operatorname{span}\{\psi_i:i\in\mathcal I\}$, and write $u=v+w$ with
$v=Pu=\sum_{\mathcal I}x_i\psi_i$ and $w=(I-P)u$. The plane waves diagonalize
$a_\theta$, so $a_\theta$ and the unweighted $L^2$ product have no cross terms
between $v$ and $w$. On the complement $a_\theta[w]\ge\Lambda_{\mathcal I}\|w\|^2$
and $b[w]\le\rho_{\max,s}\|w\|^2$, hence $(a_\theta-\lambda b)[w]\ge
c_{\mathcal I}\|w\|^2$. The cross term is
$-2\lambda\Re\langle\rho_sv,w\rangle=-2\lambda\Re\langle(I-P)\rho_sv,w\rangle$, and
$\|(I-P)\rho_sv\|^2=x^*G_{\mathcal I}x$. Completing the square,
\[
  (a_\theta-\lambda b)[u]\ge x^*T^-_{\mathcal I}x
  +c_{\mathcal I}\Bigl\|w-\frac{\lambda}{c_{\mathcal I}}(I-P)\rho_sv\Bigr\|^2 .
\]
For $w=0$ we get $(a_\theta-\lambda b)[v]=x^*T_{\mathcal I}x$, and the min--max
principle gives the left inequality. For the right one, let $V$ be a subspace on
which $a_\theta-\lambda b$ is negative definite, of dimension
$\#\{j:\lambda_{s,j}(\theta)<\lambda\}$. If $u\in V\setminus\{0\}$ had $Pu=0$, then
$(a_\theta-\lambda b)[u]\ge c_{\mathcal I}\|u\|^2>0$, a contradiction; so $P$ is
injective on $V$, and by the display $T^-_{\mathcal I}$ is negative definite on
the coefficient vectors of $PV$. Hence $n(T^-_{\mathcal I})\ge\dim V$.

\subsection{Proofs of \texorpdfstring{\cref{lem:uc,prop:reduction}}{Lemma 3.1 and Proposition 3.2}}\label{sm:proof-reduction}
\emph{\Cref{lem:uc}.} Take the side $+$. The rods do not meet $\Gamma_+$, so for
some $\epsilon>0$ the coefficient equals $\rho_{\mathrm e}$ in
$(x_+-\epsilon,x_++\epsilon)\times I_d$. Extend $v$ by zero to
$(x_+-\epsilon,x_+)\times I_d$ and quasi-periodically in $y$. Since both traces
vanish, the extension is a distributional solution of
$-\Delta\tilde v=k^2\rho_{\mathrm e}\tilde v$ in
$(x_+-\epsilon,x_++\epsilon)\times\RR$, hence real analytic; it vanishes on an
open set and therefore in the whole band. The background part of $B_+$ is
connected, so $v$ vanishes there by analytic continuation. The outer Cauchy
data on every rod boundary then vanish, so do the inner data by the
transmission conditions, and Green's representation formula inside the rod
gives $v=0$.

\emph{\Cref{prop:reduction}.} If $u\in\ker(A(\beta)-k^2)$, then
$u|_{\mathcal C_0}\in\mathcal U_0(k)$ and $u|_{B_\pm}\in\mathcal V_\pm(k)$; since
$u\in H^1(B)$ and $\Delta u\in L^2(B)$, both traces along $\nu_\pm$ agree from
the two sides, so $\mathscr T_\pm(u|_{\mathcal C_0})\in\mathcal G_\pm(k)$.
If $u|_{\mathcal C_0}=0$, then $\mathscr T_\pm(u|_{B_\pm})=0$ and
\cref{lem:uc} gives $u=0$. Conversely, for
$u_0\in\mathcal U_0^{\mathcal G}(k)$ choose $v_\pm\in\mathcal V_\pm(k)$ with
$\mathscr T_\pm v_\pm=\mathscr T_\pm u_0$ (unique by \cref{lem:uc}) and glue
$v_-,u_0,v_+$. Equal value traces give $u\in H^1_\beta(B)$; equal normal traces
mean that $\Delta u$ has no singular part on $\Gamma_\pm$, so
$\Delta u=-k^2\rho u\in L^2(B)$. Hence $u\in D(A(\beta))$ and $A(\beta)u=k^2u$.

\subsection{Proof of \texorpdfstring{\cref{thm:center-rep}}{Theorem 4.2}}\label{sm:proof-center-rep}
On $\partial\Omega_0$ the superscripts $+$ and $-$ denote traces from outside and
inside, and $\nu$ points outward.

\emph{Rayleigh coefficients.} Since $\overline{\Omega_0}\subset\mathcal C_0$,
there are background strips between $\Gamma_-$ and $\Omega_0$ and between
$\Omega_0$ and $\Gamma_+$, where $u$ is a Rayleigh series. Let $f^\pm_m,g^\pm_m$ be
the coefficients of $u|_{\Gamma_\pm}$ and $\partial_xu|_{\Gamma_\pm}$ with respect
to $\psi_m$. Define $a^-_{0,m}=\frac12(f^-_m+g^-_m/(\ii\gamma_m))$ (right-going in
the left strip) and $b^+_{0,m}=\frac12(f^+_m-g^+_m/(\ii\gamma_m))$ (left-going in the
right strip); these are unique because $\gamma_m\ne0$. With
$\langle\beta_m\rangle=(1+\beta_m^2)^{1/2}$, $u\in H^1$ near the ports gives
$\sum\langle\beta_m\rangle|f^\pm_m|^2<\infty$ and
$\sum\langle\beta_m\rangle^{-1}|g^\pm_m|^2<\infty$. By \cref{ass:wood} and
$|\gamma_m|\sim|\beta_m|$ there is $c>0$ with $|\gamma_m|\ge c\langle\beta_m\rangle$,
so $\sum\langle\beta_m\rangle(|a^-_{0,m}|^2+|b^+_{0,m}|^2)<\infty$. The waves
$r^{x_-}_m$ ($x\ge x_-$) and $l^{x_+}_m$ ($x\le x_+$) do not grow, evanescent ones
satisfy $\Im\gamma_m\ge c'\langle\beta_m\rangle$, and so their $H^1(\mathcal C_0)$
norms squared are $O(\langle\beta_m\rangle)$. Different $m$ are orthogonal in $y$
on every vertical line, hence $u_{\mathrm{hom}}$ converges in $H^1(\mathcal C_0)$.
If $\Omega_0=\emptyset$, $\mathcal C_0$ is one background strip and
$u=u_{\mathrm{hom}}$. Let now $\Omega_0\ne\emptyset$.

\emph{Auxiliary problem.} Let $w^+=u-u_{\mathrm{hom}}$ in
$\mathcal C_0\setminus\overline{\Omega_0}$ and $w^-=u$ in $\Omega_0$. By construction
$w^+$ is outgoing at both ports. Consider the whole-plane problem with exchanged
wavenumbers ($\kappa_{\mathrm e}$ in $\Omega_0$, $\kappa_{\mathrm c}$ outside with the
Sommerfeld condition) and jumps $v^+-v^-=-w^+-w^-$,
$\partial_\nu v^+-\partial_\nu v^-=-\partial_\nu w^+-\partial_\nu w^-$. By assumption it
has a unique solution $v$. Define
\[
  \tau:=w^+-v^-=-w^--v^+,\qquad
  \omega:=-\partial_\nu w^++\partial_\nu v^-=\partial_\nu w^-+\partial_\nu v^+ ,
\]
the two expressions being equal by the jump conditions.

\emph{Interior representation.} For $\mathbf z\in\Omega_0$, Green's formula for
$w^-$ gives $\mathcal S^{(\kappa_{\mathrm c})}[\partial_\nu w^-]-
\mathcal D^{(\kappa_{\mathrm c})}[w^-]=w^-$, and the exterior Green formula for the
outgoing $\kappa_{\mathrm c}$-solution $v^+$ gives
$\mathcal S^{(\kappa_{\mathrm c})}[\partial_\nu v^+]-\mathcal D^{(\kappa_{\mathrm c})}[v^+]=0$
inside. Adding, with the second expressions for $\tau,\omega$,
$\mathcal S^{(\kappa_{\mathrm c})}[\omega]+\mathcal D^{(\kappa_{\mathrm c})}[\tau]=w^-$.

\emph{Exterior representation.} Apply Green's second identity to $w^+$ and
$G^{(\kappa_{\mathrm e})}_\QP(\mathbf z,\cdot)$ on $\mathcal C_0\setminus\overline{\Omega_0}$.
The top and bottom contributions cancel by quasi-periodicity; at each port both
functions are outgoing mode by mode, so the modal Wronskians vanish. Hence
$-\mathcal S^{(\kappa_{\mathrm e})}_\QP[\partial_\nu w^+]+\mathcal D^{(\kappa_{\mathrm e})}_\QP[w^+]=w^+$.
Since $v^-$ solves the $\kappa_{\mathrm e}$ equation in $\Omega_0$ and
$G_\QP(\mathbf z,\cdot)$ is smooth there for $\mathbf z\notin\Omega_0$ (no image of
$\mathbf z$ lies in $\Omega_0$),
$\mathcal S^{(\kappa_{\mathrm e})}_\QP[\partial_\nu v^-]-\mathcal D^{(\kappa_{\mathrm e})}_\QP[v^-]=0$.
Adding, with the first expressions,
$\mathcal S^{(\kappa_{\mathrm e})}_\QP[\omega]+\mathcal D^{(\kappa_{\mathrm e})}_\QP[\tau]=w^+$.

\emph{Uniqueness.} If two sets of data give the same field, the port extraction
shows that the Rayleigh coefficients agree. For the density differences $\tau,\omega$
let $L_{\mathrm e}=\mathcal S^{(\kappa_{\mathrm e})}_\QP[\omega]+\mathcal D^{(\kappa_{\mathrm e})}_\QP[\tau]$
and $L_{\mathrm c}=\mathcal S^{(\kappa_{\mathrm c})}[\omega]+\mathcal D^{(\kappa_{\mathrm c})}[\tau]$;
$L_{\mathrm e}$ vanishes outside and $L_{\mathrm c}$ inside the rod. Since $G_\QP-G$
is smooth near $\partial\Omega_0$, the jump relations give $L_{\mathrm e}^-=-\tau$,
$\partial_\nu L_{\mathrm e}^-=\omega$, $L_{\mathrm c}^+=\tau$,
$\partial_\nu L_{\mathrm c}^+=-\omega$. Then $v^-=L_{\mathrm e}$ in $\Omega_0$ and
$v^+=-L_{\mathrm c}$ outside solve the exchanged problem with zero jumps, so $v=0$
and $\tau=\omega=0$. This is the complementary-field argument of
\cite[App.~A]{BarnettGreengard2010} (going back to \cite{ColtonKress1983}); the
port terms were removed in the first step.

\emph{Uniqueness for the exchanged problem.} If the jumps vanish, Green's first
identity in $\Omega_0$ shows that $\int_{\partial\Omega_0}\overline{v^-}\partial_\nu v^-$ is
real, hence so is $\int_{\partial\Omega_0}\overline{v^+}\partial_\nu v^+$; on
$B_R\setminus\overline{\Omega_0}$ this gives
$\Im\int_{\partial B_R}\overline v\,\partial_rv=0$, and Rellich's lemma
\cite{ColtonKress1983} gives $v^+=0$. Then $v^-$ has zero Cauchy data and vanishes
by unique continuation.

\begin{corollary}[near-image split]\label{sm:cor-near-image}
Let \cref{ass:wood} hold and $\kappa=\kappa_{\mathrm e}$.
(i) $G_{\mathrm{far}}$ in \eqref{eq:near-image} is real analytic on
$E=\RR^2\setminus\{nd\,e_y:|n|\ge2\}$ and solves $(\Delta+\kappa^2)G_{\mathrm{far}}=0$
there. (ii) Let $\overline\Omega\subset\RR\times I_d$,
$\tau\in H^{1/2}(\partial\Omega)$, $\omega\in H^{-1/2}(\partial\Omega)$,
$\delta=d/2-\max_{\partial\Omega}|y|>0$ and $S_\delta=\RR\times(-d/2-\delta,d/2+\delta)$.
Then, away from $\partial\Omega$ and $\partial\Omega\pm de_y$ in $S_\delta$,
\[
  \mathcal S_\QP[\omega]+\mathcal D_\QP[\tau]=\sum_{n=-1}^{1}q^n
  \bigl(\mathcal S[\omega]+\mathcal D[\tau]\bigr)(\cdot-nd\,e_y)+u_{\mathrm{corr}},
\]
where $u_{\mathrm{corr}}(\mathbf z)=\int_{\partial\Omega}(G_{\mathrm{far}}(\mathbf z-\mathbf z')
\omega(\mathbf z')+\partial_{\nu'}[G_{\mathrm{far}}(\mathbf z-\mathbf z')]\tau(\mathbf z'))\,\dd s'$
is smooth on all of $S_\delta$ and solves the Helmholtz equation. (iii) Under the
hypotheses of \cref{thm:center-rep}, $u$ in $\mathcal C_0\setminus\overline{\Omega_0}$
is the sum of the near-image potentials, $u_{\mathrm{corr}}$, and $u_{\mathrm{hom}}$,
where $u_{\mathrm{corr}}$ extends to a Helmholtz solution near the closed cell.
\end{corollary}

\begin{proof}
(i) In the sense of distributions $(\Delta+\kappa^2)G_\QP(\cdot,\mathbf 0)=
-\sum_nq^n\delta_{nde_y}$ \cite{Linton1998}, and $(\Delta+\kappa^2)G(\cdot,nde_y)=
-\delta_{nde_y}$; the difference solves the homogeneous equation on $E$ and is real
analytic by interior regularity. (ii) Insert \eqref{eq:near-image} into the layer
potentials. For $\mathbf z\in S_\delta$ and $\mathbf z'\in\partial\Omega$,
$|y-y'|<d$, so $\mathbf z-\mathbf z'$ stays at distance more than $d$ from the
images $|n|\ge2$; the kernel is jointly smooth there, and differentiation under the
duality pairing is allowed. (iii) Apply (ii) to \cref{thm:center-rep}.
\end{proof}

Thus the exact center field has the structure of the trial field
\eqref{eq:mfs-field}: the proxies need only approximate a solution that is smooth
near the closed cell. No convergence theory for this approximation is used; its
error is measured by the jumps.

\subsection{Piecewise regularity of the trial field}\label{sm:proof-piecewise-h2}
\begin{proposition}\label{sm:prop-piecewise-h2}
Assume that the image curves and proxies are at positive distance from the closed
cells, that the densities are trigonometric polynomials, and that
\eqref{eq:mfs-field} uses one $\kappa_{\mathrm e}$ in the background. Then
$\widehat u$ is $H^2$ up to the boundary in each homogeneous piece and has ordinary
value and normal traces on both sides of every wall and seam.
\end{proposition}

\begin{proof}
Far images, proxies and plane waves are smooth near the closed pieces, so it
suffices to consider $v=\mathcal S^{(\kappa)}[\omega]+\mathcal D^{(\kappa)}[\tau]$ on
one curve $\Gamma$. In the sense of distributions $(\Delta+\kappa^2)v=
-\omega\delta_\Gamma+\operatorname{div}(\tau\nu\delta_\Gamma)$. In normal coordinates
let $w=0$ inside and $w=(\tau(s)-t\omega(s))\chi(t)$ outside, with a smooth cutoff
$\chi$. Then $w$ has value jump $\tau$ and normal jump $-\omega$, so
$(\Delta+\kappa^2)w=F-\omega\delta_\Gamma+\operatorname{div}(\tau\nu\delta_\Gamma)$ with
$F\in L^2_{\mathrm{loc}}$. Hence $(\Delta+\kappa^2)(v-w)=-F$ and $v-w\in
H^2_{\mathrm{loc}}$, while $w$ is smooth up to $\Gamma$ from either side. Near walls
and seams the one-sided cell formulas extend smoothly across, so there are no
corner singularities.
\end{proof}

\subsection{Proof of \texorpdfstring{\cref{prop:trace}}{Proposition 5.1}}\label{sm:proof-trace}
We use arc length $s$ on $\Gamma$, the normal map
$\boldsymbol\Phi_\Gamma(s,t)=\boldsymbol r_\Gamma(s)+t\nu_\Gamma(s)$, and the signed
curvature $\kappa_\Gamma$ with $\boldsymbol r_\Gamma''=\kappa_\Gamma\nu_\Gamma$, so that
$\nu_\Gamma'=-\kappa_\Gamma\boldsymbol r_\Gamma'$ and the area element is
$(1-t\kappa_\Gamma(s))\,\dd s\,\dd t$. For a convex rod with outward normal,
$\kappa_\Gamma<0$; walls have $\kappa_\Gamma=0$. The condition $h_g<1$ keeps the
Jacobian positive; injectivity excludes nonlocal self-intersections (necks) and is
checked separately (\cref{sm:reliable-cg}). Collars that cross the top or bottom
edge are understood on the quasi-periodic cylinder.

\begin{lemma}[global $H^2$ regularity]\label{sm:global-h2}
For $v\in D(A(\beta))$, $v\in H^2(B)$ and
$\|D^2v\|_{L^2(B)}=\|\Delta v\|_{L^2(B)}\le\rax^{1/2}\|A(\beta)v\|_{\mathcal H}$. In
particular, the value and normal traces of $v$ agree from both sides of every
interface.
\end{lemma}

\begin{proof}
Extend $v$ quasi-periodically in $y$ and expand in $\psi_m$ in $y$ and by the
Fourier transform in $x$. By Parseval,
$\sum_m\int(\xi^2+\beta_m^2)^2|\widehat v_m(\xi)|^2\,\dd\xi=\|\Delta v\|^2<\infty$, and
the same sum equals $\|D^2v\|^2$. Finally $\Delta v=-\rho A(\beta)v$.
\end{proof}

\begin{proof}[Proof of \cref{prop:trace}]
Fix one normal segment, let $\delta=\delta_g$, $\kappa=\kappa_\Gamma(s)$ and
$w(t)=1-\kappa t$, so $\int_{-\delta}^{\delta}w=2\delta$. Averaging
$g(0)=g(t)-\int_0^tg'$ against $w(t)\,\dd t$ and exchanging the integrals gives
$g(0)=\frac1{2\delta}\int wg+\int K_\kappa g'$ with
\[
  |K_\kappa(t)|=\begin{cases}
  \bigl(\delta-t-\kappa(\delta^2-t^2)/2\bigr)/(2\delta),&0<t<\delta,\\
  \bigl(\delta+t-\kappa(t^2-\delta^2)/2\bigr)/(2\delta),&-\delta<t<0.
  \end{cases}
\]
Weighted Cauchy--Schwarz gives $|g(0)|\le(2\delta)^{-1/2}(\int w|g|^2)^{1/2}
+T_\kappa^{1/2}(\int w|g'|^2)^{1/2}$ with $T_\kappa=\int|K_\kappa|^2/w$, and a direct
computation with $h=|\kappa|\delta$ gives
$T_\kappa=\frac\delta6+\frac{\delta h^2}8\int_0^1\frac{(1-x)^4}{1-h^2x^2}\,\dd x\le\tau_g$.
Apply $(a+b)^2\le(1+\eta)a^2+(1+\eta^{-1})b^2$ and integrate in $s$; since
$w\,\dd s\,\dd t$ is the area element, summing over the group costs a factor $N_g$.
Using this for $g=v\circ\boldsymbol\Phi_\Gamma$ ($|g'|\le|\nabla v|$) and for
$g=(\nu_\Gamma\cdot\nabla v)\circ\boldsymbol\Phi_\Gamma$ ($|g'|\le|D^2v|$) gives the
two trace inequalities. Finally, with the spectral measure $E_\lambda$ of
$A(\beta)$, supported in $[\lambda_0,\infty)$, $\|v\|^2_{L^2}\le\rin^{-1}\|v\|^2_{\mathcal H}$,
$\|\nabla v\|^2=(A(\beta)v,v)_{\mathcal H}$, and \cref{sm:global-h2},
$T_g(v)^2\le\int p_g(\lambda)\,\dd\|E_\lambda v\|^2\le C_g^2\int(\lambda+\mu)^2\,
\dd\|E_\lambda v\|^2=C_g^2\|(A(\beta)+\mu)v\|^2_{\mathcal H}$.
\end{proof}
